%% Lebenslauf - Ana Lorena Ortiz Loyola
%% Zugeschnitten auf: Bonpago GmbH - Werkstudent (m/w/d) für angewandte KI & KI-Agenten, Frankfurt am Main
%% Deutsche Version von main_bonpago.tex.
%%
%% Kompilieren mit: lualatex main_bonpago_de.tex
%% (pdflatex schlägt auf modernem MiKTeX häufig mit fontawesome5-Fehlern fehl.)

\documentclass[11pt,a4paper,sans]{moderncv}
\moderncvstyle{banking}
\moderncvcolor{blue}

\usepackage{needspace}

% Erzwingt, dass Vor- und Nachname sowie die Sektionsüberschriften in moderncv
% Blau (color1) dargestellt werden. Der banking-Standardstil rendert diese auf
% lualatex+MiKTeX sonst schwarz, was inkonsistent mit dem restlichen blauen
% Akzentschema wirkt.
\renewcommand*{\firstnamestyle}[1]{{\fontsize{34}{36}\bfseries\upshape\color{color1}#1}}
\renewcommand*{\lastnamestyle}[1]{{\fontsize{34}{36}\bfseries\upshape\color{color1}#1}}
\renewcommand*{\sectionstyle}[1]{{\sectionfont\color{color1}#1}}

\usepackage[utf8]{inputenc}
\usepackage{hyperref}
\hypersetup{
    colorlinks=true,
    linkcolor=blue,
    filecolor=magenta,
    urlcolor=blue,
    pdftitle={Ana Lorena Ortiz Loyola - Lebenslauf},
    pdfpagemode=FullScreen,
}
\usepackage[scale=0.80]{geometry}
\usepackage{import}

% Persönliche Daten
\name{Ana Lorena}{Ortiz Loyola}
\address{Frankfurt am Main, Deutschland}{}{}
\phone[mobile]{+49 176 62145797}
\email{a01750283@tec.mx}
\extrainfo{\href{https://linkedin.com/in/lorenasolana}{LinkedIn}}

\begin{document}

\makecvtitle

% ============================================================
%     PROFILTEXT
% ============================================================

\vspace{6pt}
\small{Data-Science-Studentin (Tecnológico de Monterrey, Abschluss voraussichtlich 2027) mit praktischer Erfahrung im Einsatz eines KI-gestützten Programmierassistenten (Claude Code) zur Umsetzung eines realen End-to-End-Projekts sowie großem Interesse daran, dies professionell auf die Evaluierung und Entwicklung von KI-Agenten und agentischen Workflows auszuweiten. Kombiniert eine fundierte Machine-Learning-Grundlage (Scikit-Learn, Keras, TensorFlow) mit praktischer Data-Engineering-Erfahrung bei der Deutschen Bundesbank sowie der Fähigkeit, technische Arbeit klar zu dokumentieren und zu kommunizieren.}

% ============================================================
%     KERNKOMPETENZEN
% ============================================================

\section{Kernkompetenzen}
\vspace{1pt}
\begin{itemize}

\item \textbf{Angewandte KI-Tools}: Praktische Erfahrung mit einem KI-gestützten Programmierassistenten (Claude Code) in einem realen Projekt; großes Interesse an der Evaluierung von KI-Tools, Modellen und agentischen Workflows

\item \textbf{Machine Learning}: Scikit-Learn, Keras, TensorFlow, Random Forest, SVM, K-Means

\item \textbf{Programmierung}: Python, R, SQL, Java, C++

\item \textbf{Data Engineering}: ETL/ELT-Pipelines, Automatisierung der Datenqualität, Migrationstests (Deutsche Bundesbank AnaCredit)

\item \textbf{Dokumentation \& Kommunikation}: Strukturierte Berichterstattung, Stakeholder-Kommunikation, MS Office (Word, Excel, PowerPoint)

\item \textbf{Cloud \& Tools}: Azure Cognitive Services, AWS, Git

\end{itemize}

% ============================================================
%     BERUFSERFAHRUNG
% ============================================================

\section{Berufserfahrung}
\vspace{3pt}
\begin{itemize}

\needspace{5\baselineskip}
\item{\cventry{04/2026--07/2026}{Data Engineer / Praktikantin}{Deutsche Bundesbank (AnaCredit)}{Frankfurt am Main, Deutschland}{}{\vspace{1pt}
\begin{itemize}
    \item Entwickelte automatisierte Migrationstests (Zeilenanzahl, referenzielle Integrität, Aggregatkonsistenz) mit Pandas-DataFrames zur Validierung der Geschäftslogik nach der Datenmigration
    \item Entwickelte automatisierte Fehlererkennung zur Identifikation statistischer Abweichungen in großen gemeldeten Bankdaten, inklusive automatisiertem Jira-Ticketing und Stakeholder-Kommunikation
\end{itemize}}}

\end{itemize}

% ============================================================
%     PROJEKTE
% ============================================================

\needspace{8\baselineskip}
\section{Projekte}
\vspace{1pt}
\begin{itemize}

\item{\cventry{05/2026}{Banking-Datenpipeline -- Datenbereinigung \& Prüfprotokoll}{Persönliches Projekt (GitHub), umgesetzt mit Claude Code}{}{}{\vspace{1pt}
\begin{itemize}
    \item Konzipierte und entwickelte eine End-to-End-Datenpipeline für simulierte Banktransaktionen (150 Konten) unter Einsatz von Claude Code als KI-gestütztem Programmierassistenten zur Beschleunigung der Umsetzung
    \item Entwickelte eine normalisierte SQLite-Datenbank mit Fremdschlüssel-Integritätsprüfungen, mehrstufiger Datenbereinigung und einem vollständig nachvollziehbaren Protokoll jeder Änderung
    \item Erstellte Qualitätsvisualisierungen und einen strukturierten Ausnahmebericht
\end{itemize}}}

\vspace{3pt}

\item{\cventry{10/2024--12/2024}{Predictive Modeling -- Optimierung im Gesundheitswesen}{Tecnológico de Monterrey}{}{}{\vspace{1pt}
\begin{itemize}
    \item Entwickelte und validierte Predictive-Modelle (Random Forest, neuronale Netze) sowie K-Means-Clustering auf Basis von über 25 Jahren nationaler Gesundheitsdaten
\end{itemize}}}

\vspace{3pt}

\item{\cventry{04/2023--06/2023}{Luftqualitätsanalyse Mexiko-Stadt -- SVM-Klassifikation}{Tecnológico de Monterrey}{}{}{\vspace{1pt}
\begin{itemize}
    \item Entwickelte SVM-Modelle (linearer und RBF-Kernel) zur Klassifikation kritischer Ozonwerte anhand stadtweiter NOx/O3-Daten
\end{itemize}}}

\end{itemize}

% ============================================================
%     AUSBILDUNG
% ============================================================

\section{Ausbildung}
\vspace{1pt}
\begin{itemize}

\item{\cventry{2022--2027}{BSc in Data Science}{Tecnológico de Monterrey}{Mexiko}{}{\vspace{1pt}
Note: 1,3 (Stipendium für akademische Exzellenz; DAAD-KOSPIE-Stipendium). Schwerpunkte: Machine Learning \& KI, Datenanalyse \& Statistik, numerische Optimierung.
}}

\vspace{3pt}

\item{\cventry{2025--2026}{Auslandssemester}{Universität Göttingen}{Deutschland}{}{}}

\end{itemize}

% ============================================================
%     SPRACHEN
% ============================================================

\section{Sprachen}
\vspace{1pt}
\begin{itemize}
\item Deutsch (C1), Englisch (C1, TOEFL iBT 105/120), Spanisch (Muttersprache).
\end{itemize}

% ============================================================
%     PUBLIKATIONEN
% ============================================================

\section{Publikationen}
\vspace{1pt}
\begin{itemize}
\item Beitrag zu einer wissenschaftlichen Publikation über gitterbasierte Kryptografie (Sicherheitsparameter- und Laufzeitanalyse von LWE und NTRU), gemeinsam mit Prof. A. F. De Abiega L'Eglisse, Tecnológico de Monterrey (Zitation ausstehend, Publikation in Vorbereitung).
\end{itemize}

% ============================================================
%     AUSZEICHNUNGEN
% ============================================================

\section{Auszeichnungen}
\vspace{1pt}
\begin{itemize}
\item{\textbf{Stipendium für akademische Exzellenz} -- Tecnológico de Monterrey.}
\item{\textbf{DAAD-KOSPIE-Stipendium}.}
\item{\textbf{Nationale Physik-Olympiade} -- Mexiko (2020).}
\item{\textbf{Lobende Erwähnung, Physik-Olympiade} (2021).}
\end{itemize}

% ============================================================
%     REFERENZEN
% ============================================================

\section{Referenzen}
\vspace{1pt}
\begin{itemize}
\item{Auf Anfrage erhältlich.}
\end{itemize}

\end{document}
